\documentclass[../main.tex]{subfiles}
\begin{document}

\chapter{Conclusiones y trabajo futuro}
\label{C:conclusiones}

\section{Conclusiones}

Este trabajo ha presentado el diseño, implementación y evaluación de un laboratorio
de simulación de ransomware basado en contenedores Docker con un \emph{stack}
\gls{siem} completo. El trabajo ha cubierto el ciclo completo de un proyecto de
investigación experimental en ciberseguridad: revisión bibliográfica, diseño del
entorno, implementación, ejecución de 120~experimentos automatizados, análisis
estadístico riguroso y validación con técnicas reales de la matriz MITRE ATT\&CK.

\subsection{Cumplimiento de los objetivos específicos}

Los siete objetivos específicos definidos en la introducción han sido alcanzados:

\begin{description}
  \item[OE1 — Laboratorio reproducible] Se ha diseñado e implementado un
  laboratorio de nueve contenedores Docker completamente automatizado, publicado
  como software libre en GitHub, que puede reproducirse en cualquier máquina Linux
  con Docker instalado en aproximadamente diez minutos. La validación en máquina
  virtual Ubuntu~24.04 confirmó esta reproducibilidad.

  \item[OE2 — Tres estrategias de backup] Se han implementado y comparado las
  tres estrategias definidas: control negativo sin backup, repositorio \restic
  sobre SFTP convencional y repositorio \restic sobre MinIO con Object Lock
  COMPLIANCE. Las diferencias entre estrategias han quedado completamente
  cuantificadas.

  \item[OE3 — Cuatro escenarios de ataque] Se han implementado y validado los
  cuatro escenarios correspondientes a las técnicas T1486, T1490, T1021.004 y
  T1041 de la matriz MITRE ATT\&CK.

  \item[OE4 — Instrumentación con Wazuh] Se han desarrollado 13~reglas de
  detección custom que cubren los cuatro escenarios mediante \gls{fim} y eventos
  \gls{json} estructurados, adaptadas a las limitaciones técnicas del entorno.

  \item[OE5 — Diseño factorial y análisis estadístico] Se han ejecutado las
  120~ejecuciones del diseño factorial y se han analizado los resultados con
  Mann-Whitney~U, delta de Cliff, \emph{bootstrap} de intervalos de confianza y
  test exacto de Fisher.

  \item[OE6 — Validación con Atomic Red Team] Las cuatro técnicas ATT\&CK han
  sido validadas con tests de Atomic Red Team ejecutados sobre el laboratorio
  real, confirmando que las reglas detectan técnicas de ataque reales más allá
  de los escenarios sintéticos.

  \item[OE7 — Publicación como software libre] El repositorio completo está
  disponible públicamente en GitHub con documentación suficiente para la
  reproducción independiente de los experimentos.
\end{description}

\subsection{Sobre la hipótesis H1}

La comparación estadística del \gls{mttr} entre las estrategias~B y~C no muestra
diferencias significativas en los escenarios donde ambas pueden recuperarse
(E1: $p = 0.473$; E3: $p = 0.623$; E4: $p = 0.186$). Los tamaños del efecto son
pequeños o negligibles en E1 y E3, y solo E4 alcanza un efecto mediano
(Cliff~$\delta = 0.360$) sin alcanzar significación estadística.

\textbf{La hipótesis H1 se rechaza}: la ventaja del almacenamiento inmutable no
reside en la velocidad de restauración, sino en la garantía de recuperabilidad
cuando el repositorio ha sido destruido.

\subsection{Sobre la hipótesis H2}

Los tiempos de detección mediante \gls{siem} oscilan entre 5.7 y 41.9~segundos,
frente a los más de 200~días que la literatura reporta para la detección manual.
La diferencia de más de cinco órdenes de magnitud es concluyente.

\textbf{La hipótesis H2 se confirma} con amplísimo margen: la instrumentación con
un \gls{siem} transforma cualitativamente la capacidad de respuesta ante incidentes
de ransomware, pasando de una detección que puede tardar meses a una detección que
tarda segundos.

\subsection{Sobre la hipótesis H3}

El test exacto de Fisher sobre el escenario~E2 produce $p < 0.001$ (0/10 éxitos
para B vs. 10/10 para C). La magnitud de la diferencia es la máxima posible.

\textbf{La hipótesis H3 se confirma}: Object Lock COMPLIANCE proporciona una
ventaja \emph{cualitativa} sobre el \emph{backup} convencional al garantizar la
recuperabilidad ante la destrucción activa del repositorio.

\subsection{Sobre el laboratorio como contribución}

El laboratorio desarrollado constituye en sí mismo una contribución al campo. Sus
características principales son:

\begin{itemize}
  \item \textbf{Reproducibilidad total}: cualquier investigador puede replicar el
  experimento clonando el repositorio.
  \item \textbf{Automatización completa}: las 120~ejecuciones se completan en
  cuatro horas sin intervención manual.
  \item \textbf{Extensibilidad modular}: nuevos escenarios de ataque o estrategias
  de backup pueden añadirse sin modificar la estructura central.
  \item \textbf{Metodología estadística rigurosa}: pruebas no paramétricas,
  tamaños del efecto e intervalos de confianza \emph{bootstrap}.
  \item \textbf{Validación con herramientas reales}: Atomic Red Team confirma que
  las reglas detectan técnicas reales más allá de los escenarios del experimento.
\end{itemize}

\subsection{Reflexión final}

La respuesta que proporciona este trabajo es inequívoca: una estrategia de
\emph{backup} convencional, por bien implementada que esté, es vulnerable a un
atacante que obtiene las credenciales del servidor de \emph{backup}. El
almacenamiento inmutable con Object Lock COMPLIANCE elimina esa vulnerabilidad de
forma efectiva y con un coste adicional modesto respecto a una solución SFTP
convencional.

La combinación de almacenamiento inmutable con un \gls{siem} bien configurado
proporciona una defensa en profundidad que eleva significativamente el coste y la
complejidad del ataque para el adversario: no solo garantiza que habrá algo que
restaurar, sino que la organización sabrá que está siendo atacada en cuestión de
segundos.

\section{Trabajo futuro}

Las siguientes líneas de trabajo futuro se identifican como extensiones naturales
de este estudio:

\begin{enumerate}[label=TF\arabic*., leftmargin=2.5cm]

  \item \textbf{Aumento del tamaño de muestra}: ampliar las repeticiones por celda
  a al menos 30 para detectar efectos medianos con potencia del 80\,\%. Esto
  permitiría confirmar o refutar la diferencia marginal de E4
  (Cliff~$\delta = 0.360$, $p = 0.186$) y proporcionar intervalos de confianza
  más estrechos.

  \item \textbf{Cifradores reales en entorno sandbox}: sustituir el cifrador
  sintético de Python por muestras reales de \emph{LockBit}~3.0 o \emph{BlackCat}
  en un entorno aislado sin acceso exterior. Esto produciría tiempos de cifrado
  más representativos e identificaría diferencias en los patrones de acceso a
  ficheros que afecten al \gls{fim}.

  \item \textbf{Ciclo de vida completo del ataque}: implementar las fases de
  reconocimiento previo al ataque y la persistencia del agente malicioso para
  simular el comportamiento de un \gls{apt} real. Esto permitiría evaluar la
  detección temprana antes del cifrado.

  \item \textbf{Evaluación con múltiples snapshots en estrategia B}: implementar
  una política de retención con \emph{snapshots} horarios y evaluar cuántos
  \emph{snapshots} necesita el atacante destruir antes de que el \gls{siem} active
  procedimientos de contención.

  \item \textbf{Respuesta activa automática de Wazuh}: integrar las capacidades
  de \emph{active response} de \wazuh (aislamiento de red, bloqueo de procesos)
  y medir el impacto sobre el daño total del ataque.

  \item \textbf{Validación con Atomic Red Team para T1490 en Linux}: extender la
  validación cuando estén disponibles tests oficiales de Atomic Red Team para Linux
  en la técnica T1490.

  \item \textbf{Detección basada en aprendizaje automático}: entrenar modelos de
  detección de anomalías (Isolation Forest, Autoencoder) sobre los datos de
  \emph{log} generados. La hipótesis es que estos modelos podrían detectar el
  ataque antes de que los umbrales de frecuencia del \gls{fim} se superen.

  \item \textbf{Medición del overhead de monitorización FIM}: cuantificar
  sistemáticamente el impacto en \gls{cpu} y E/S de disco que introduce el
  \gls{fim} en tiempo real sobre el \emph{fileserver} durante el cifrado activo,
  para determinar si la monitorización introduce algún sesgo en los tiempos medidos
  y si es viable en servidores con alta carga de trabajo.

  \item \textbf{Comparación con el modo Governance de Object Lock}: este trabajo
  evalúa exclusivamente el modo COMPLIANCE con 30~días de retención. No se comparó
  con el modo Governance, que permite a usuarios con permisos especiales anular la
  protección, ni con períodos de retención distintos (7, 14 o 90~días). Un análisis
  de la relación entre período de retención, coste de almacenamiento y nivel de
  protección aportaría criterios prácticos para adaptar la estrategia~C a distintos
  presupuestos y requisitos normativos.

  \item \textbf{Paneles de visualización en el dashboard de Wazuh}: desarrollar
  paneles y gráficos en el dashboard de Wazuh que permitan a un operario de
  seguridad visualizar en tiempo real la evolución de las alertas: distribución
  por escenario, nivel de alerta, agente origen y línea temporal de detección.
  Esta visualización facilitaría la monitorización continua sin necesidad de
  consultar la \gls{api} REST directamente.

  \item \textbf{Barrido sistemático de umbrales de frecuencia \gls{fim}}:
  como se plantea en la sección~\ref{sec:sensibilidad} del
  capítulo~\ref{C:discusion}, realizar un análisis de sensibilidad que ejecute
  el escenario~E1 con distintas combinaciones de \texttt{frequency} y
  \texttt{timeframe} en las reglas~100200 y~100202, midiendo el compromiso
  resultante entre tiempo de detección y tasa de falsos positivos ante
  actividad legítima de copia masiva de ficheros. Esto permitiría sustituir la
  calibración actual, basada en observación durante la fase de desarrollo, por
  una curva de compromiso caracterizada empíricamente.

  \item \textbf{Extensión a orquestadores de contenedores distintos de Docker
  Compose}: evaluar la portabilidad de la arquitectura descrita en el
  capítulo~\ref{C:arquitectura} a Kubernetes, sustituyendo las comprobaciones
  de \texttt{healthcheck} de Docker Compose (sección~\ref{sec:recursos-healthchecks})
  por sondas de disponibilidad (\emph{readiness}/\emph{liveness probes})
  nativas. Esta extensión resultaría especialmente relevante para organizaciones
  que ya operan sus cargas de trabajo sobre Kubernetes y que deseasen adoptar
  el laboratorio sin introducir una segunda herramienta de orquestación de
  contenedores en su infraestructura.

\end{enumerate}

\section{Difusión y aplicabilidad como recurso docente}
\label{sec:difusion}

Más allá de su contribución como trabajo de investigación experimental, el
diseño modular del laboratorio --- ya destacado como característica principal
en la sección anterior --- le confiere un valor potencial como recurso docente
para la formación en ciberseguridad defensiva. La separación entre la
infraestructura de orquestación (\texttt{scripts/}), la lógica de cada
escenario de ataque (\texttt{simulators/}) y las reglas de detección
(\texttt{lab-config/wazuh/}), descrita en el anexo~C, permite emplear
subconjuntos del laboratorio de forma aislada: un curso introductorio a
\gls{fim} podría limitarse al escenario~E1 y a las reglas~100200--100206, sin
necesidad de desplegar el resto del \emph{stack} de detección de movimiento
lateral o exfiltración.

Esta aplicabilidad docente es coherente con la tendencia, documentada en el
capítulo~\ref{C:estado_arte} a través del trabajo de Bitzki et
al.~\cite{netsecbed2026} sobre \emph{testbeds} nativos en contenedores, de
emplear infraestructuras ligeras y reproducibles tanto para investigación como
para formación práctica en ciberseguridad, en contraposición a los laboratorios
tradicionales basados en máquinas virtuales completas, cuyo coste de
despliegue y mantenimiento dificulta su uso en un contexto docente con
recursos limitados. El requisito de recursos documentado en el anexo~C
(16\,GB de \gls{ram}, cuatro núcleos) es compatible con el equipamiento
informático habitual de un aula universitaria de nivel de grado.

La licencia MIT bajo la que se publica el código, descrita en la última
sección del anexo~C, permite explícitamente esta reutilización con fines
docentes, incluyendo la modificación de los escenarios existentes o la
incorporación de escenarios adicionales por parte de terceros, sin más
condición que mantener el aviso de \emph{copyright} original.

\section{Consideraciones éticas}

Todo el trabajo se ha realizado en un entorno completamente aislado de la red
pública: los contenedores operan en una red privada Docker sin conexión a Internet
durante la ejecución de los escenarios de ataque. Los \emph{scripts} de ataque
desarrollados (el cifrador, la destrucción de repositorios, el movimiento lateral
y la exfiltración) son herramientas de propósito específico diseñadas para el
laboratorio y carecen de los mecanismos de propagación, evasión de antivirus y
comunicación con servidores de mando y control que caracterizarían a un
\emph{malware} real.

El repositorio público del proyecto incluye únicamente el código del laboratorio y
la infraestructura de automatización, no código malicioso reutilizable fuera del
contexto del laboratorio. El cifrador incluido genera una clave aleatoria para
cada ejecución que solo existe en memoria durante esa ejecución, garantizando que
los ficheros cifrados durante el experimento no pueden descifrase sin la clave de
esa ejecución específica.

La validación con Atomic Red Team se realizó usando herramientas del propio sistema
objetivo siguiendo las directrices de uso responsable del \emph{framework}, que
establece que los tests deben ejecutarse únicamente sobre sistemas propios o con
autorización explícita.

\end{document}
